% --- IMPORTS ---
\documentclass[leqno]{article}
\usepackage{graphicx}
\usepackage{dsfont}
\usepackage{mathtools}
\usepackage{amssymb}
\usepackage{amsmath}
\usepackage{amsthm}
\usepackage{float}
\usepackage{ragged2e}
\usepackage{tensor}
\usepackage{fancyhdr}
\usepackage{empheq}
\usepackage[most]{tcolorbox}
\usepackage{enumitem}
\usepackage{dirtytalk}
\usepackage{marginnote}
\usepackage{microtype}
\usepackage[
  a4paper,
  left=0.8in,
  right=1in,
  top=1in,
  bottom=0.5in,
  headheight=14pt,
  headsep=0.4in
]{geometry}
\setlength{\parindent}{0cm}
% --- TikZ Packages ---
\usepackage{pgfplots}
\pgfplotsset{compat=1.18}
\usepgfplotslibrary{fillbetween, polar}
\usetikzlibrary{calc, positioning}
\usetikzlibrary{angles, quotes}
\usetikzlibrary{arrows.meta}
\usetikzlibrary{decorations.pathreplacing, decorations.markings}
\usetikzlibrary{patterns}
\usetikzlibrary{intersections}
% --- URL Packages ---
\usepackage[colorlinks=true,linkcolor=red,citecolor=magenta,urlcolor=black]{hyperref}%
\urlstyle{same}
% --- END OF IMPORTS ---

% --- CUSTOM COMMANDS ---
\RenewDocumentCommand{\marks}{o m}{%
  \IfValueTF{#1}%
    {\marginnote{\raggedright\textbf{[#2]}}[#1]}%
    {\marginnote{\raggedright\textbf{[#2]}}}%
}
\newcommand{\vspacerule}[2]{%
  \par\vspace*{#1}%
  \noindent\hrulefill\par
  \vspace*{#2}%
}
\definecolor{scarlet}{RGB}{205, 40, 75}
\newcommand{\questionitem}{%
  \item\phantomsection
  \addcontentsline{toc}{section}{Question \number\value{enumi}}%
}
\renewcommand{\contentsname}{Questions}
% --- END OF CUSTOM COMMANDS ---

% --- HEADER CONFIG ---
\pagestyle{fancy}
\fancyhead{}
\fancyfoot{}
\renewcommand{\headrulewidth}{0.9pt}
\lhead{Vectors - Shortest Distances}
\chead{}
\rhead{\href{https://danieljsmith.org}{danieljsmith.org}}
% --- END OF HEADER CONFIG ---

\begin{document}
\vspace*{20pt}
\tableofcontents
\newpage
\begin{enumerate}[
  label=\textbf{\arabic*.},
  leftmargin=*,
  topsep=0pt,
  partopsep=0pt,
  parsep=0pt
]

% ---- Question 1 ----
\questionitem

The plane $\Pi$ has equation $3x+y+2z=-4$. The point $P$ has coordinates $(5,-5,3)$. \\
\begin{enumerate}[label=\textbf{(\alph*)}]
  \item Calculate the shortest distance from $P$ to $\Pi$. \marks{2} \\
\end{enumerate}
The line $L$ has vector equation
\[
\mathbf{r}=\begin{pmatrix}2\\1\\-3\end{pmatrix}+\lambda\begin{pmatrix}1\\-2\\2\end{pmatrix}
\] \\
\begin{enumerate}[label=\textbf{(\alph*)}, resume]
  \item Verify that $P$ lies on $L$. \marks{2} \\
  \item Find the coordinates of the point where $L$ meets $\Pi$. \marks{3} \\
  \item Determine the acute angle between $L$ and $\Pi$. \marks{4} \\
  \item Use the results of parts (b), (c) and (d) to verify your answer to part (a). \marks{3} \\
\end{enumerate}

\newpage

% ---- Question 2 ----
\questionitem

The points $C$ and $D$ have coordinates $(0,0,2)$ and $(3,3,-1)$ respectively.\\

The planes $P_1$ and $P_2$ have equations $x+y+z=3$ and $2x-y=1$ respectively. \\
\begin{enumerate}[label=\textbf{(\alph*)}]
  \item Find the acute angle between the line $CD$ and the plane $P_1$. \marks{4} \\
  \item Show that the line $CD$ meets $P_1$ and $P_2$ in a common point, and find its coordinates. \marks{5} \\
  \item
  \begin{enumerate}[label=(\roman*)]
    \item Find $(\mathbf{i}+\mathbf{j}+\mathbf{k})\times(2\mathbf{i}-\mathbf{j})$. \marks{1}
    \item Hence find the acute angle between the planes $P_1$ and $P_2$. \marks{3}
    \item Find the shortest distance between the point $C$ and the line of intersection of the planes\\ $P_1$ and $P_2$. \marks{4} \\
  \end{enumerate}
\end{enumerate}

\newpage

% ---- Question 3 ----
\questionitem

The plane $\Pi$ has equation
\[
\mathbf{r}=\begin{pmatrix}2\\-1\\3\end{pmatrix}+\lambda\begin{pmatrix}1\\2\\1\end{pmatrix}+\mu\begin{pmatrix}4\\-1\\-2\end{pmatrix}
\]
where $\lambda$ and $\mu$ are scalar parameters. \\
\begin{enumerate}[label=\textbf{(\alph*)}]
  \item Show that vector $\mathbf{i}-2\mathbf{j}+3\mathbf{k}$ is perpendicular to $\Pi$. \marks{2} \\
  \item Hence find a Cartesian equation of $\Pi$. \marks{2} \\
\end{enumerate}
The line $\ell$ has equation
\[
\mathbf{r}=\begin{pmatrix}-2\\1\\4\end{pmatrix}+t\begin{pmatrix}1\\1\\1\end{pmatrix}
\]
where $t$ is a scalar parameter. \\

The point $A$ lies on $\ell$. Given that the shortest distance between $A$ and $\Pi$ is $\tfrac{9}{\sqrt{14}}$ \\
\begin{enumerate}[label=\textbf{(\alph*)}, resume]
  \item determine the possible coordinates of $A$. \marks{4} \\
\end{enumerate}

\newpage

% ---- Question 4 ----
\questionitem

The line $\ell_1$ has vector equation
\[
\mathbf{r}=\begin{pmatrix}2\\-1\\3\end{pmatrix}+\lambda\begin{pmatrix}1\\2\\-1\end{pmatrix}
\] \\
The line $\ell_2$ is the line of intersection of the planes
\[
x-y+z=6
\] \\
and
\[
2x+y-z=3
\] \\
\begin{enumerate}[label=(\roman*)]
  \item Find the shortest distance between $\ell_1$ and $\ell_2$. \marks{5} \\
  \item Find a Cartesian equation of the plane which contains $\ell_1$ and is parallel to $\ell_2$. \marks{2} \\
\end{enumerate}

\newpage

% ---- Question 5 ----
\questionitem

\begin{enumerate}[label=\textbf{(\alph*)}]
  \item Show that the three planes with equations
  \[
  \begin{aligned}
  x+\lambda y+z&=4 \\
  2x-y+2z&=3 \\
  2x+y-z&=5
  \end{aligned}
  \]
  where $\lambda$ is a constant, meet at a unique point except for one value of $\lambda$, which is to be determined. \marks[-20pt]{3}\\
  \item In the case $\lambda=2$, use matrices to find the point of intersection $P$ of the planes, showing your method clearly. \marks{3} \\
\end{enumerate}
The line $\ell$ has equation
\[
\frac{x+1}{1}=\frac{y-1}{-1}=\frac{z}{1}
\] \\
\begin{enumerate}[label=\textbf{(\alph*)}, resume]
  \item Find a vector equation of $\ell$. \marks{2} \\
  \item Find the shortest distance between the point $P$ and $\ell$. \marks{4} \\
  \item
  \begin{enumerate}[label=(\roman*)]
    \item Show that $\ell$ is parallel to the plane $2x+y-z=5$. \marks{3}
    \item Find the distance between $\ell$ and the plane $2x+y-z=5$. \marks{2} \\
  \end{enumerate}
\end{enumerate}

\newpage

% ---- Question 6 ----
\questionitem

\begin{enumerate}[label=\textbf{(\alph*)}]
  \item The plane $\Pi$ passes through the points $A(1,-2,3)$, $B(2,0,3)$ and $C(2,-2,2)$.\\
  
  The point $D$ has coordinates $(4,1,0)$. Find the shortest distance between $D$ and $\Pi$. \marks{4} \\
  \item Lines $\ell_1$ and $\ell_2$ have the following equations. \\
  \[
  \ell_1:\ \mathbf{r}=\begin{pmatrix}1\\2\\-1\end{pmatrix}+\lambda\begin{pmatrix}2\\1\\2\end{pmatrix}
  \] \\
  \[
  \ell_2:\ \mathbf{r}=\begin{pmatrix}3\\-1\\4\end{pmatrix}+\mu\begin{pmatrix}1\\-2\\1\end{pmatrix}
  \] \\
  Find, in exact form, the distance between $\ell_1$ and $\ell_2$. \marks{5} \\
\end{enumerate}

\newpage

% ---- Question 7 ----
\questionitem

The position vectors of the points $A$, $B$, $C$, $D$ are
\[
\mathbf{i}+2\mathbf{j},\quad 2\mathbf{i}+2\mathbf{j},\quad \mathbf{i}+3\mathbf{j}+2\mathbf{k},\quad 2\mathbf{i}+3\mathbf{j}+m\mathbf{k}
\]
respectively, where $m$ is an integer. \\

It is given that the shortest distance between the line through $A$ and $D$ and the line through $B$ and $C$ is $\tfrac{1}{\sqrt{14}}$ \\
\begin{enumerate}[label=\textbf{(\alph*)}]
  \item Show that the only possible value of $m$ is $1$. \marks{7} \\
  \item Find the shortest distance of $C$ from the line through $A$ and $D$. \marks{3} \\
  \item Show that the acute angle between the planes $ABD$ and $ACD$ is $\cos^{-1}\left(\frac{\sqrt{3}}{2}\right)$. \marks{4} \\
\end{enumerate}

\newpage

% ---- Question 8 ----
\questionitem

\begin{enumerate}[label=\textbf{(\alph*)}]
  \item Given that $\mathbf{u}=\lambda\mathbf{i}-2\mathbf{j}+2\mathbf{k}$ and $\mathbf{v}=\mathbf{i}+\mathbf{j}-\mathbf{k}$, find the following, giving your answers in terms of $\lambda$. 
  \begin{enumerate}[label=(\roman*)]
    \item $\mathbf{u}\cdot\mathbf{v}$ \marks{1}
    \item $\mathbf{u}\times\mathbf{v}$ \marks{2} \\
  \end{enumerate}
  \item Hence determine
  \begin{enumerate}[label=(\roman*)]
    \item the acute angle between the planes $2x-2y+2z=11$ and $x+y-z=7$ \marks{3}
    \item the shortest distance between the lines
    \[
    \mathbf{r}=\mathbf{i}-\mathbf{k}+s(\mathbf{i}-2\mathbf{j}+2\mathbf{k})
    \] 
    and
    \[
    \mathbf{r}=2\mathbf{i}+3\mathbf{j}+3\mathbf{k}+t(\mathbf{i}+\mathbf{j}-\mathbf{k})
    \]\\
    giving your answer as a multiple of $\sqrt{2}$. \marks{3}
  \end{enumerate}
\end{enumerate}
\end{enumerate}
\end{document}
